\section{Discussion}
\label{sec:discussion}

In this national, leakage-audited evaluation of 368,132 isolated-CABG operations across 1,091 U.S.\ surgical centers, we resolve a central tension in perioperative clinical machine learning: why published POAF algorithms routinely report apparent AUCs above 0.85 yet collapse to a median external AUC of 0.61 in critical appraisals,\cite{fields2025poaf,lopeztuesta2026poaf} and how disciplined validation, continuous-target learning, parallel genetic programming, and hierarchical institutional calibration together establish a trustworthy clinical benchmark.

\subsection{Why Validation Geometry Matters as Much as Model Architecture}
Our forensic audit of STS ACSD v2.81 demonstrates that apparent performance in national surgical registries can be distorted by three silent mechanisms before a single hyperparameter is tuned. First, ignoring parent--child skip logic (\cli{Complics == 1}) discards 63.3\% of uncomplicated recoveries and induces textbook collider bias.\cite{hernan2020causal} Second, positional column slicing retains seven early-column postoperative duration variables that masquerade as preoperative inputs. Third, default final-year temporal holdouts (calendar 2017) retain 99.5\% hospital overlap while mixing half-year coverage with missingness drift, hiding a $-1.99$-AUC-point external drop when identical models are evaluated on strictly unseen hospitals.

Our controlled comparison between strictly unseen hospitals (0\% overlap) and seen hospitals (100\% all-1,091-hospitals-in-train) quantifies this distortion directly. Merely placing patients from every hospital into training inflated patient-only discrimination by $+28$ to $+64$ bps across every model family even with hospital identifiers completely withheld, because flexible learners implicitly fingerprint site-level coding, laboratory-ordering, and medication-prescribing patterns ($r=0.088$--$0.143$ with true hospital POAF incidence). Frontier tabular foundation models---TabPFN-3.5, KumoTabular-Large v2, and TabFM---gained nearly twice as much from hospital familiarity ($+49$ to $+64$ bps) as Dual-Paradigm Fusion ($+28$ bps), demonstrating that in-context tabular transformers should never be evaluated solely on hospital-overlapping splits when cold-start multi-center transport is the intended clinical use case.

\subsection{Breaking the Binary Ceiling and Reconciling Cold-Start vs.\ In-Network Deployment}
Under strictly unseen-hospital external validation on 207 raw pre-incision predictors, all six world-leading 50k binary learning families compressed within a narrow 35-bp band (0.6712--0.6746), reflecting the intrinsic biological ceiling of pre-incision registry data given unmeasured intraoperative trauma, atrial stretch, pericardial inflammation, and autonomic surges.\cite{echahidi2008mechanisms,dobrev2019postoperative} Both of our proposed paradigms broke above that ceiling without test leakage: 1,000-generation evolutionary continuous-target distillation (EvoTarget-FT) stabilized without collapse and lifted all seven model families---including $+27$ to $+97$-bp gains on zero-shot tabular foundation model regressors---while our six-fix parallel genetic programming engine (JAIN Ultra GP v2) delivered a $+39.8$-bp rescue over legacy JAIN. Their orthogonal fusion achieved 0.6767, decisively surpassing every bedside score ($+316$ to $+1{,}029$ bps) and every 50k binary AI baseline ($+20$ to $+55$ bps, plus $+9.9$ bps over a 3-way binary blend) while recovering 79.2\% of the 7$\times$ national data-scaling gap.

Crucially, our dual-regime evaluation clarifies how cardiothoracic AI should actually be deployed in clinical practice. For \textbf{cold-start deployment at a new hospital} (zero historical local registry data), Dual-Paradigm Fusion provides the strongest discrimination (0.6767), near-ideal national calibration (slope~0.983, O:E~0.994, ECE~0.0072), transparent physiological attribution ($R^2 = 0.89$, including independently protective preoperative amiodarone), and consistent equity across 24 of 25 demographic, comorbidity, and volume subgroups. Once a center joins an established national registry or multi-hospital health system (\textbf{in-network deployment}), our seen-hospital study proves that combining calibrated patient-level predictions with a strictly historical Empirical-Bayes hospital incidence prior yields a massive $+174$ to $+213$-bp leap---bringing 50k models to 0.696--0.699 and 346k CatBoost past 0.70 to 0.7014, while cutting prospective 2017 center-level rate error by 14\% with just a single local intercept update. Furthermore, our observation that lower-entanglement 10k-context foundation models pair slightly better with explicit 346k hospital intercepts than higher-entanglement 40k--346k models motivates future explicitly hierarchical mixed-effects neural and symbolic training that separates patient-level pathophysiology from institutional intercepts during fitting.

\subsection{Clinical Implications for Targeted POAF Prophylaxis}
In cardiothoracic surgery, discrimination matters only insofar as it improves bedside decisions. Dual-Paradigm Fusion separated strictly unseen CABG patients by 7.4-fold in observed POAF risk---from 6.1\% in D1 (where aggressive prophylaxis risks net harm) to 44.8\% in D10 (where nearly half of patients develop treatment-requiring AF). In Decision Curve Analysis across the clinically relevant 25\%--35\% decision zone for interventions such as posterior left pericardiotomy (PALACS\cite{gaudino2021palacs}), perioperative amiodarone loading, or colchicine prophylaxis, Dual-Paradigm Fusion generated strictly positive Net Benefit and avoided 20.4 to 33.8 unnecessary interventions per 100 operations compared with universal prophylaxis, outperforming the widely cited Mariscalco POAF score in 100\% of hospital-cluster bootstrap draws.

\subsection{Limitations}
Five limitations warrant explicit transparency. First, although all test labels remained cryptographically isolated during training and inference, our iterative methodological benchmark evaluated more than 140 model configurations on the same 74-hospital holdout; accordingly, final confirmatory deployment requires freezing the exact Dual-Paradigm Fusion and hierarchical site-intercept pipeline prior to prospective testing on newer STS ACSD versions (v4.20+). Second, STS \cli{COtAFib} captures clinically recognized, treatment-requiring postoperative AF or flutter rather than asymptomatic brief episodes detected by continuous research telemetry.\cite{filardo2017telemetry,perezgrovas2023jaha} Third, STS ACSD v2.81 lacks raw preoperative 12-lead ECG waveforms, echocardiographic left atrial volume index, and circulating inflammatory or natriuretic biomarkers, which likely explain the residual discrimination ceiling near 0.68 (cold-start) to 0.70 (in-network).\cite{han2025aiecg,khalpey2026digital} Fourth, while our strictly nested 207-raw EvoTarget-FT control confirmed $+12.3$ bps of zero-overlap denoising (65\% of total gain), its 74-hospital cluster CI spanned zero ($[-3.3, +28.5]$, $P>0=0.942$), reflecting finite cluster power across 74 centers. Fifth, institutional incidence priors improve in-network discrimination by capturing center-level ascertainment and care-pathway differences alongside case mix, and therefore must be monitored longitudinally for practice drift.

\subsection{Conclusions}
In a rigorously audited national cohort of 368,132 isolated-CABG procedures across 1,091 U.S.\ centers, combining continuous-target clinical self-attention (EvoTarget-FT) with re-engineered parallel genetic programming (JAIN Ultra GP v2) decisively outperforms established bedside scores and every frontier binary tabular model across strictly unseen hospitals, delivers clinically actionable Net Benefit and demographic equity, and---when paired with empirical-Bayes institutional calibration in seen hospitals---establishes a principled, dual-regime foundation for national perioperative AF risk stratification.
